% Layout macros for the portrait. They only place ordinary notation on the page:
% numerals, + and \cdot, [ ], floor, ceiling and definite integrals.
% One font, one weight (New Computer Modern); every delimiter is a .5bp stroke.
\usepackage{amsmath}
\usepackage{unicode-math}
\usepackage{graphicx}
\setmainfont{NewCM10-Book.otf}
\setmathfont{NewCMMath-Book.otf}[StylisticSet=2]   % upright integral signs
\newcommand{\Num}{\fontsize{9.4bp}{9.4bp}\selectfont}
\newcommand{\Scr}{\fontsize{5.6bp}{5.6bp}\selectfont}
\newcommand{\At}[3]{\hbox to0pt{\kern\dimexpr#1\relax\raise\dimexpr#2\relax\hbox{#3}\hss}}
% --- within one line of the expression (lengths in bp, heights from the line's baseline)
\def\K#1{\kern#1bp}                                               % move right
\newcommand{\Lift}[2]{\raise#1bp\hbox{#2}}                         % a numeral at mid-height of a tall bracket
\newcommand{\Plus}[2]{\raise#2bp\hbox to#1bp{\hss$\mathord{+}$\hss}}        % the operator +
\newcommand{\Times}[2]{\raise#2bp\hbox to#1bp{\hss$\mathord{\cdot}$\hss}}   % the operator \cdot
% delimiters drawn from height #1 to height #2
\newcommand{\Stroke}[2]{\vrule width.5bp height#2bp depth-#1bp}
\newcommand{\HookB}[1]{\vrule width.45bp height\dimexpr#1bp+.45bp\relax depth-#1bp}
\newcommand{\HookT}[1]{\vrule width.45bp height#1bp depth-\dimexpr#1bp-.45bp\relax}
\newcommand{\OpenBracket}[2]{\Stroke{#1}{#2}\HookB{#1}\kern-.45bp\HookT{#2}\kern-.45bp}    % [
\newcommand{\CloseBracket}[2]{\kern-.45bp\HookB{#1}\kern-.45bp\HookT{#2}\Stroke{#1}{#2}}   % ]
\newcommand{\OpenFloor}[2]{\Stroke{#1}{#2}\HookB{#1}\kern-.45bp}                           % \lfloor
\newcommand{\CloseFloor}[2]{\kern-.45bp\HookB{#1}\Stroke{#1}{#2}}                          % \rfloor
\newcommand{\OpenCeil}[2]{\Stroke{#1}{#2}\HookT{#2}\kern-.45bp}                            % \lceil
\newcommand{\CloseCeil}[2]{\kern-.45bp\HookT{#2}\Stroke{#1}{#2}}                           % \rceil
% \Integral{y}{height}{a}{b}{f}{lift}{x}  =  \int_a^b f \, dx  with a tall sign:
% the hooks are the font's integral pieces, the stem between them a rule of the same thickness
\newbox\Ia \newbox\Ib
\newcommand{\Integral}[7]{\hbox to0pt{\kern-2.39bp\fontsize{3.57bp}{3.57bp}\selectfont
  \setbox\Ia\hbox{$\inttop$}\setbox\Ib\hbox{$\intbottom$}%
  \rlap{\raise\dimexpr#1bp+\dp\Ib\relax\copy\Ib}%
  \rlap{\raise\dimexpr#1bp+#2bp-\ht\Ia\relax\copy\Ia}%
  \kern2.39bp\vrule width.5bp height\dimexpr#1bp+#2bp-\ht\Ia+.1bp\relax depth-\dimexpr#1bp+\ht\Ib+\dp\Ib-.1bp\relax\hss}%
  \hbox to0pt{\kern1.5bp\raise\dimexpr#1bp-.1bp\relax\hbox{\Scr#3}\hss}%
  \hbox to0pt{\kern2.9bp\raise\dimexpr#1bp+#2bp-4.6bp\relax\hbox{\Scr#4}\hss}%
  \hbox to0pt{\kern2.6bp\raise#6bp\hbox{$#5\,\mathrm{d}#7$}\hss}}
% --- a sum whose lines are stacked inside one tall pair of brackets:  \BeginSum{width}{height} lines \EndSum
\newcommand{\BeginSum}[2]{\hbox to#1bp\bgroup\vrule width0pt height#2bp depth0pt
  \At{.3bp}{.6bp}{\vrule width.5bp height\dimexpr#2bp-1.2bp\relax depth0pt}%
  \At{.8bp}{.6bp}{\vrule width.6bp height.45bp depth0pt}\At{.8bp}{#2bp-1.05bp}{\vrule width.6bp height.45bp depth0pt}%
  \At{#1bp-.8bp}{.6bp}{\vrule width.5bp height\dimexpr#2bp-1.2bp\relax depth0pt}%
  \At{#1bp-1.4bp}{.6bp}{\vrule width.6bp height.45bp depth0pt}\At{#1bp-1.4bp}{#2bp-1.05bp}{\vrule width.6bp height.45bp depth0pt}}
\newcommand{\EndSum}{\hss\egroup}
\newcommand{\Row}[1]{\hbox to0pt\bgroup\raise#1bp\hbox\bgroup}   % one line, baseline #1 above the floor of the sum
\newcommand{\EndRow}{\egroup\hss\egroup}
\newcommand{\Factor}{\hbox to0pt\bgroup\lower.3bp\hbox\bgroup}   % a bracketed sum used as a factor in a line
\newcommand{\EndFactor}{\egroup\hss\egroup}
